% Be sure to set \appendix or \appendices in structural/body.tex
\chapter{Literature review}

There is no question that these sorts of orbital games have become an ever more prevalent area of research in aerospace autonomy academia, with plenty of examples of related works in recent years, mostly on optimal-control-based methods.

In \textit{“Pursuit--Evasion Game for Satellites Based on Continuous Thrust Reachable Domain”}, \cite{gong2020continuousthrust}, the authors developed a classical differential game of satellite pursuit-evasion under continuous thrust using a Clohessy--Wiltshire (CW) model. Moreover, \textit{“A Pursuit-Evasion Game in the Orbital Plane”}, \cite{selvakumar2017orbitalplane}, they develop a two player zero-sum differential game with input constraints, providing a semi-analytical capture-time solution to their problem using closed-form optimal controls. Then in \textit{“Multiple-to-One Orbital Pursuit: A Computational Game Strategy”}, \cite{liu2025multipletoone}, they propose a one-step look-ahead Nash equilibrium method for cooperative pursuit-evasion in cooperative spacecraft. And lastly, in \textit{"Low-Thrust Orbital Differential Games with Speed Constraint Enforcement Using Cost Weighting"}, \cite{drucker2026lowthrustgames}, they frame a low-thrust spacecraft pursuit-evasion problem as a linear-quadratic zero-sum game with terminal speed constraints, from which they obtain guidance laws via cost-weight selection demonstrating improved performance against actively evading targets.


In addition to more classical approaches to these rigorous space-bound problems, there’s been a rise of journal published research trying to solve variations of these games but opting for online reinforcement learning (RL) approaches in recent years. Despite the long standing criticism for core reinforcement learning algorithms to not be able to provide theoretical convergence guarantees to learned policies, researchers have found problems in which they've come in useful and provided empirically convergent policies, strong performance relative to their assigned tasks, and excellent computational efficiency.


In \textit{"Online Solution for Orbital Pursuit-Evasion Game via Heterogeneous Proximal Policy Optimization"}, \cite{li2025heterogeneousppo}, they successfully developed a PPO framework for multi-agent orbital pursuit-evasion, outperforming shooting methods in their optimality index metrics in computational efficiency while providing robustness to suboptimal opponent strategies. Next, in \textit{A Non-Fixed-Duration Turn-Based Orbital Game Algorithm Based on History-Aware and Predictive Result Reward Proximal Policy Optimization}, \cite{cao2026historyawareppo}, they proposed HPRR-PPO algorithm to do alternating learning for non-fixed-duration turn-based orbital game problems, showing effective interception against evaders with compute speed suitable for real time deployment claiming a compute runtime of 1 ms per made decision. Then in \textit{"Near-Optimal Interception Strategy for Orbital Pursuit-Evasion using Deep Reinforcement Learning"}, \cite{zhang2022nearoptimal}, they use deep reinforcement learning to achieve near-optimal guidance laws in one to one orbital pursuit-evasion games. Finally, in \textit{"{PRD-MADDPG}: An Efficient Learning-Based Algorithm for Orbital Pursuit-Evasion Game with Impulsive Maneuvers"}, \cite{zhao2023prdmaddpg} propose a MADDPG based method for impulsive orbital pursuit-evasion.

The fact that all these papers above have been published in the past few years makes it clear
that there is growing interest in developing reinforcement learning-based methods for planning in adversarial orbital dynamics. This is especially relevant given the generalized growing interest within the autonomy community to implement machine learning-based frameworks to complement or even replace classical control architectures. 


Furthermore, additional work has been made in the areas of delivering estimation-aware RL frameworks. While physical AI is something roboticists and industry alike are pushing towards, one must remember that real-life robots won't have full-state measurements of their environment as we can achieve in a simulation. Instead, they need to use sensors to characterize their surroundings and plan accordingly. Generally speaking, this is problematic on the decision-making side, as it can turn a clean Markov Decision Process (MDP) that you'd see in a full-state environment into a Partially Observable Markov Decision Process (POMDP), which is harder to characterize (\cite{kaelbling1998pomdp}). In “DESPOT: Online POMDP Planning with Regularization”, \cite{ye2017despot}, the authors show that modern decision-making under partial observability can be carried out by planning over an approximated belief tree, showing that near-optimal policies can be obtained in partial observability without solving for the full POMDP. In "Deep Variational Reinforcement Learning for POMDPs", \cite{igl2018dvrl}, the authors show that reinforcement learning under partial observability using latent-state inference to estimate, meaning that a belief representation is learned from incomplete observations, showing that this method outperforms baselines in partially observable tasks. In "Learning Deep Neural Network Policies with Continuous Memory States", \cite{zhang2015continuousmemory}, they do partially observable learning with continuous control, where they handle the problem by augmenting the policies with continuous memory states, where they end up finding that learned policies can memorize relevant information needed to act with incomplete observations.

Moreover, some works have tried to fill in the gaps between learning frameworks and classical control frameworks by combining the most desirable parts of each into a single integrated architecture.

In "Safe Model-based Reinforcement Learning with Stability Guarantees", \cite{berkenkamp2017safembrl} , they show that learning can be constrained via classical stability/safety filtering, using Lyapunov safety certificates instead of leaving the learning unconstrained. In "Learning-based Model Predictive Control for Safe Exploration and Reinforcement Learning", \cite{koller2019safeexplorationrl} they bridge safe explorations and reinforcement learning through model predictive control, showing that uncertainty-aware trajectories and terminal safety constraints can be made safely while improving learning performance And most importantly for the sake of this work, in “Safe Reinforcement Learning Using Robust Control Barrier Functions” ,\cite{emam2021saferlrcbf}, they show how a Control-Barrier-Function (CBF) safety layer can be used to modify the learned policy's behavior by projecting the policy actions in a way that induces a safe learning transitions map, preserving the safety of the learned model while allowing a reward-based policy to be learned on the side.

We notice a few gaps in the existing literature. Namely, we see that most current orbital RL work, including those works cited in this literature review, assume privileged, full-state information within their frameworks, and practically none of the presented works study learned policies with online estimation in the loop. Almost none of these papers study the interaction of RL policies when intertwined with safety filters and classical control in the loop. 

All the above justify the novelty of this work. Here, we are not inventing any novel fundamental RL algorithm or the sort, but rather the novelty we present in this work is the integration of these pieces into a coherent and stable PARL+BC framework. Here we're using a combination of reinforcement learning techniques that can lead to coherent, robust learning strategies in zero-sum orbital defense game settings: using PPO to find learned policies under the given set of hyperparameters and complex reward shaping, doing so without help of an equilibrium/Nash prior solution, and testing how to implement this method under a partial-observability setting via bearings-only Kalman Filter measurements. In this work, we contribute to the present literature in the following ways:

\begin{itemize}
    \item Using a staged PPO training pipeline instead of residual networks to train the agents' policies in a zero-sum, two player orbital game. Namely, the defender and attacker policies are trained in alternating phases against frozen and randomized mixtures of opponent policies. The goal of this is to produce policies that become more specialized, yet more generalizable, to complete their given task within the game framework instead of overfitting to any one rival. 

    \item Extending upon the PPO defense framework from the previous bullet to partially observable measurements via behavior cloning. In other words, this work evaluates policies where an agent doesn't have access to the states of their adversaries. Instead, these have to be measured via bearings-only Kalman-filter estimates, based on the work by \cite{aidala1979kalman}. This provides a controlled test on how learned orbital defense policies degrade under sensor and partial observability limitations that would be present in a real world implementation.
    \item Studying the interaction between learned policies in full and partial observability within the context of a classical CBF-QP safety filter. This shows how a model-based classical controller modifies or constrains learning behavior during training and rollouts.
\end{itemize}

So in summary, existing literature shows that orbital pursuit evasion remains a relevant GNC problem to research and that reinforcement learning is increasingly being used by various works to obtain guidance policies. Nevertheless, most of the existing learning-based work cited above assumes access to privileged state information and separates the learned guidance problem from a spacecraft's estimation and safety-based controls architecture. 

The contributions made by this work are therefore not the invention of any fundamental RL, guidance, estimation, or control algorithms (e.g., PPO, Kalman filtering, or CBF-QP control, to give a few examples), but rather the integration of all these algorithms into a phased learning framework.

Specifically, by developing the PARL+BC framework, we research how alternating, PPO-trained attacker and defender policies perform given opponent randomization, bearings-only estimator measurements, and CBF-QP safety control filtering. This positions the work done on PARL+BC as an effective architecture and case study for work on learning-based orbital defense subject to realistic constraints one would see in a real world autonomous vehicle, spacecraft or otherwise.
